\subsection{Production Release Process and Capture Reliability}
\label{sec:release-process}

The Samsung camera software studied here is a privileged preloaded system application that interacts with Samsung-specific Android framework extensions, system services, vendor components, and companion applications.
Capture is a core cross-layer workflow that spans these components, so a failure discovered late can require a firmware update rather than an application-only fix.
Because firmware updates require cross-component coordination and broad revalidation, daily development builds undergo automated black-box testing across supported camera workflows, device models, and OS versions.
Periodically, the camera-validation team and development engineers stress-test the latest builds, diagnose timing and memory-pressure failures, and verify fixes.

\smallskip\noindent\textbf{Capture Timeout.}
To maintain reliable capture across diverse devices and runtime conditions, the product treats compliance with its fixed internal Capture Timeout deadline as a high-priority camera reliability KPI and a hard release gate.
For each capture, the timeout clock starts when the framework accepts the application's request and stops once the required frames and metadata are collected and the \emph{Draft Sequence} produces an early user-visible \emph{Draft image}.
Each violation triggers a fail-fast application crash to expose the failure and blocks release until developers correct its cause or restrict the affected feature under those conditions.
Because a targeted restriction ships faster than a cross-component fix, such restrictions accumulate across models and features.
The deployed camera software therefore carries several at once, each disabling a capture-side capability on a different runtime proxy: the thermal level, the number of captures already in flight, or whether a mode enables a particular processing stage.
Every one of these proxies stands in for the same question, whether a capture can still meet its deadline, and none of them measures it.
